
\documentclass[12pt,a4paper]{article}
\usepackage{ctex}
\usepackage{geometry}
\geometry{margin=2.5cm}
\usepackage{listings}
\usepackage{xcolor}
\usepackage{amsmath}
\usepackage{booktabs}
\usepackage{caption}
\usepackage{enumitem}

% 全局代码块样式
\lstset{
    basicstyle=\ttfamily\small,
    frame=single,
    numbers=left,
    numberstyle=\tiny\color{gray},
    backgroundcolor=\color{gray!5},
    breaklines=true,
    columns=flexible,
    captionpos=b,
    aboveskip=1em,
    belowskip=1em
}

% 列表间距
\setlist{itemsep=0.2em, topsep=0.5em}

\begin{document}

\title{系统开发工具基础 实验报告}
\author{孔博\\25020007055}
\date{2026年8月30日}
\maketitle

\section{实验一：含空格文件名的批量整理}
\subsection{实验代码}
\begin{lstlisting}[language=bash, caption=目录创建与测试文件生成]
mkdir -p q01 && cd q01

mkdir -p input/docs input/tmp

cat > input/docs/notes.one.txt <<EOF
alpha
beta
EOF
echo "hidden" > input/docs/.secret.txt

touch input/tmp/empty.txt

cat > input/run.log <<EOF
2026-08-30 start
2026-08-30 stop
EOF
\end{lstlisting}

\begin{lstlisting}[language=bash, caption=文件复制、权限配置与清单生成]
mkdir -p work/2502007055

find input -type f -name "*.txt" -exec cp --parents {} work/2502007055/ \;

find work/2502007055 -type d -exec chmod 750 {} \;

find work/2502007055 -type f -exec chmod 640 {} \;

find work/2502007055 -type f -ls > inventory.txt
\end{lstlisting}

\subsection{实验结果}
\begin{enumerate}
  \item 目录与文件验证：执行 \verb|ls -la input| 可显示全部文件，包含以点开头的隐藏文件 \verb|.secret.txt|，符合长格式展示要求。
  \item 批量复制结果：\verb|work| 目录完整保留 \verb|input/docs|、\verb|input/tmp| 的层级结构，所有txt文件复制完成，文件名含空格可正常处理。
  \item 权限配置结果：所有子目录权限为750，所有普通文件权限为640。
  \item 清单文件：\verb|inventory.txt| 中逐条记录了文件的字节数、相对路径等信息。
\end{enumerate}

\section{实验二：随机访问日志统计}
\subsection{实验代码}
\begin{lstlisting}[language=bash, caption=analyze.sh 统计脚本]
#!/bin/bash

# 接收CSV文件路径参数
csv_file="$1"

# 文件不存在时，错误输出到stderr，返回非零退出码
if [ ! -f "$csv_file" ]; then
    echo "错误：文件 $csv_file 不存在" >&2
    exit 1
fi

echo "===== 5xx 访问最多的前 2 个 path ====="
# 提取5xx状态的路径，统计次数并按规则排序
awk -F',' 'NR>1 && $4 ~ /^5/{print $3}' "$csv_file" \
| sort \
| uniq -c \
| sort -k1,1nr -k2,2 \
| head -n 2

echo -e "\n===== 平均 latency_ms（保留两位小数） ====="
# 跳过表头，计算延迟平均值
awk -F',' 'NR>1 {sum += $5; cnt++} END{printf "%.2f\n", sum/cnt}' "$csv_file"
\end{lstlisting}

\subsection{实验结果}
\begin{enumerate}
  \item 正常文件运行输出：
\begin{lstlisting}
===== 5xx 访问最多的前 2 个 path =====
      3 /api/orders
      2 /api/users

===== 平均 latency_ms（保留两位小数） =====
193.75
\end{lstlisting}
  \item 异常场景验证：对不存在的文件执行脚本，错误信息输出至标准错误流；执行 \verb|echo $?| 返回退出码1，满足非零退出要求。
\end{enumerate}

\section{实验三：制造、解决并解释一次合并冲突}
\subsection{实验代码}
\begin{lstlisting}[language=bash, caption=仓库初始化与分支创建]
mkdir -p q03 && cd q03
git init

echo "mode=normal" > config.txt
git add config.txt
git commit -m "init config mode=normal"

git checkout -b feature-a
echo "mode=safe" > config.txt
git add config.txt
git commit -m "feature-a change to safe"

git checkout main
git checkout -b feature-b
echo "mode=fast" > config.txt
git add config.txt
git commit -m "feature-b change to fast"
\end{lstlisting}

\begin{lstlisting}[language=bash, caption=分支合并与冲突解决]
git checkout main
git merge feature-a

git merge feature-b

git add config.txt
git commit -m "resolve merge conflict: keep safe, add note=reviewed"
\end{lstlisting}

\subsection{实验结果}
\begin{enumerate}
  \item 冲突现象：合并feature-b时，\verb|config.txt| 出现Git冲突标记，\verb|<<<<<<< HEAD| 标记当前分支内容，\verb|>>>>>>> feature-b| 标记待合并分支内容。
  \item 冲突解决：最终文件保留 \verb|mode=safe|，新增一行 \verb|note=reviewed|，完全删除冲突标记。
  \item 历史验证：执行 \verb|git log --all --graph --decorate --oneline| 可查看分支分叉、合并的图形化提交历史；执行 \verb|git status| 显示工作区干净。
\end{enumerate}

\section{实验四：修复并构建一页技术说明}
\subsection{实验代码}
\begin{lstlisting}[language=tex, caption=report.tex 文档源码]
\documentclass{article}
\usepackage{amsmath}
\usepackage{booktabs}

\begin{document}
\title{Tool Report}
\author{孔博-25020007055}
\maketitle

\section{Result}

The formula \ref{eq:sum} shown below is a summation formula.
\begin{equation}
\label{eq:sum}
S = \sum_{i=1}^{n} x_i
\end{equation}

Table \ref{tab:data} shows two groups of experimental data.
\begin{table}[htbp]
  \centering
  \caption{Experimental Data}
  \label{tab:data}
  \begin{tabular}{cc}
    \toprule
    Group A & Group B \\
    \midrule
    10 & 20 \\
    30 & 40 \\
    \bottomrule
  \end{tabular}
\end{table}

\end{document}
\end{lstlisting}

编译命令：
\begin{lstlisting}[language=bash]
latexmk -pdf -halt-on-error report.tex
\end{lstlisting}

\subsection{实验结果}
\begin{enumerate}
  \item 成功生成 \verb|report.pdf|，文档结构完整，包含标题、作者、章节与正文。
  \item 公式自动编号，正文通过 \verb|\ref| 交叉引用显示正确数字，无 \verb|??| 错误。
  \item 表格包含标题与自动编号，正文引用正常，三线表格式规范。
\end{enumerate}

\end{document}
